% =============================================================================
% Chapter 9: Inference at Scale — ML Systems Lecture Slides (Volume II)
% =============================================================================
% IMAGES NEEDED (SVGs converted to PDF):
%   - serving-cost-inversion.pdf   - training-vs-inference.pdf
%   - serving-hierarchy.pdf        - batching-strategies.pdf
%   - kv-cache-wall.pdf            - paged-attention.pdf
%   - power-of-two-choices.pdf     - cold-start-timeline.pdf
%   - queuing-theory-llm.pdf       - serving-dimensions.pdf
% =============================================================================
\documentclass[aspectratio=169, 12pt]{beamer}
\usepackage{../../assets/beamerthememlsys}

\mlsyssetup{
  volume       = {Volume II},
  chapter      = {Chapter 9},
  logo         = {../../assets/img/logo-mlsysbook.png},
  instlogo     = {../../assets/img/logo-harvard.png},
  chaptertitle = {Inference at Scale},
}

% --- Fonts ---
\usepackage{fontspec}
\setsansfont{texgyreheros}[
  Extension=.otf,
  UprightFont=*-regular,
  BoldFont=*-bold,
  ItalicFont=*-italic,
  BoldItalicFont=*-bolditalic,
]
\setmonofont{texgyrecursor}[
  Extension=.otf,
  UprightFont=*-regular,
  BoldFont=*-bold,
  ItalicFont=*-italic,
  BoldItalicFont=*-bolditalic,
  Scale=0.85,
]

% --- Packages ---
\usepackage{booktabs}
\usepackage{amsmath}

% --- Image paths ---
\graphicspath{
  {images/}
}

% --- Chapter-specific macros ---
\newcommand{\TTFT}{\text{TTFT}}
\newcommand{\TPOT}{\text{TPOT}}

% --- Helper: safe image include ---
\newcommand{\safeimg}[2][width=\textwidth,keepaspectratio]{%
  \IfFileExists{images/#2}{\includegraphics[#1]{#2}}{%
    \IfFileExists{#2}{\includegraphics[#1]{#2}}{%
      \fbox{\parbox[c][2.5cm][c]{0.85\linewidth}{\centering\footnotesize\textcolor{midgray}{[Missing image]}}}%
    }%
  }%
}

% --- Section count for navigation ---
\setcounter{mlsystotalsections}{8}

\title{Inference at Scale}
\author{Vijay Janapa Reddi}
\institute{Harvard University}
\date{}

\begin{document}

% =============================================================================
% TITLE SLIDE
% =============================================================================
\mlsystitle{Inference at Scale}{Where ML Systems Live or Die Economically}{cover_inference_at_scale.png}

% =============================================================================
% LEARNING OBJECTIVES
% =============================================================================
\begin{frame}{Learning Objectives}
\note{
% -- LINK: What prior concept connects to this slide
Ch 8 orchestrated the fleet for training. This chapter shifts to
serving --- where ML systems generate revenue and where costs dominate.

% -- NARRATE: What to SAY while showing this slide
Walk through objectives. Emphasize: inference cost dominates training
cost by 10--1000x for production systems.

% -- ENGAGE: Specific question, prediction, or task for THIS slide
Ask: ``How much do you think it costs to serve ChatGPT for one day?''
Expected: students will guess thousands. Real answer: millions.

% -- WARN: What students will get wrong on THIS topic
Students think training is the expensive part because it gets
headlines (\$100M training runs). Serving costs that every week.

% -- FLEX: [CORE] or [OPTIONAL] + contingency
[CORE] Sets the economic motivation for the chapter.
IF SHORT: Read objectives quickly.
}

\footnotesize
\begin{enumerate}\setlength\itemsep{0pt}
  \item Quantify why \textbf{serving cost dominates training cost} using the total cost equation
  \item Design \textbf{batching strategies}: static, continuous, and feature-parallel
  \item Apply the \textbf{serving hierarchy} (request, replica, service, platform)
  \item Analyze \textbf{KV cache} memory challenges and \textbf{PagedAttention}
  \item Compare model sharding: \textbf{tensor}, \textbf{pipeline}, \textbf{expert parallelism}
  \item Explain why \textbf{power-of-two-choices} achieves exponentially better load balancing
  \item Design \textbf{autoscaling} policies accounting for GPU cold start latency
\end{enumerate}

\end{frame}

\begin{frame}{Visual Language}
\note{
% -- NARRATE: What to SAY while showing this slide
Briefly point to each card. Students know this from prior chapters.

% -- FLEX: [OPTIONAL] + contingency
[OPTIONAL] Skip if students are familiar from Ch 7--8.
}

\small
Throughout this course, colors carry meaning:

\vspace{0.3cm}
\begin{columns}[T]
  \begin{column}{0.45\textwidth}
    \begin{mlsyscard}{computestroke}
    \textbf{Blue} --- Compute / Processing\\
    {\footnotesize GPU ops, forward/backward pass, inference}
    \end{mlsyscard}
    \vspace{0.15cm}
    \begin{mlsyscard}{datastroke}
    \textbf{Green} --- Data / Memory\\
    {\footnotesize Data flow, caches, healthy paths}
    \end{mlsyscard}
  \end{column}
  \begin{column}{0.45\textwidth}
    \begin{mlsyscard}{routingstroke}
    \textbf{Orange} --- Routing / Scheduling\\
    {\footnotesize Load balancers, batch windows}
    \end{mlsyscard}
    \vspace{0.15cm}
    \begin{mlsyscard}{errorstroke}
    \textbf{Red} --- Error / Cost / Bottleneck\\
    {\footnotesize Loss, decode phase, waste}
    \end{mlsyscard}
  \end{column}
\end{columns}
\end{frame}



% =============================================================================
\section{Economics}
% =============================================================================

\begin{frame}{The \$5M Training Run That Costs \$5M Per Week to Serve}
\note{
% -- LINK: What prior concept connects to this slide
Ch 8 covered the economics of GPU clusters for training. This slide
reveals the economic inversion: serving dwarfs training cost.

% -- NARRATE: What to SAY while showing this slide
``Training GPT-4 class models costs millions --- a one-time expense.
But serving them costs that amount every week. Recommendation systems
have a 100--1000x serving-to-training cost ratio. Even LLM APIs run
10--100x.'' Point to the table and the cost inversion diagram.

% -- ENGAGE: Specific question, prediction, or task for THIS slide
Ask: ``What fraction of total ML cost do you think goes to inference?''
Expected: students guess 50\%. Real answer: 80--90\% for production systems.

% -- WARN: What students will get wrong on THIS topic
Students focus on training cost because it is visible (cloud bills,
GPU procurement). Serving cost is distributed across many smaller
requests and harder to track.

% -- FLEX: [CORE] or [OPTIONAL] + contingency
[CORE] The economic inversion is the central motivation for the chapter.
IF SHORT: Show the table and state the inversion; skip the diagram details.
}

\footnotesize
\begin{columns}[T]
  \begin{column}{0.50\textwidth}
    Training is a \textbf{one-time expense}. Serving is \textbf{perpetual}.

    \vspace{0.1cm}
    \begin{mlsyscard}{errorstroke}
    {\footnotesize Inference OpEx exceeds training CapEx by \textbf{orders of magnitude}.}
    \end{mlsyscard}

    \vspace{0.1cm}
    {\scriptsize
    \renewcommand{\arraystretch}{1.1}
    \begin{tabular}{@{}lr@{}}
      \toprule
      \textbf{Application} & \textbf{Serve/Train} \\
      \midrule
      Recommendation & 100--1000$\times$ \\
      Search ranking & 100--1000$\times$ \\
      LLM API        & 10--100$\times$ \\
      \bottomrule
    \end{tabular}
    }
  \end{column}
  \begin{column}{0.47\textwidth}
    \safeimg[width=\textwidth,height=5.0cm,keepaspectratio]{serving-cost-inversion.pdf}
  \end{column}
\end{columns}

\end{frame}

\begin{frame}{The Serving Cost Equation}
\note{
% -- LINK: What prior concept connects to this slide
The previous slide showed the inversion qualitatively. This slide
provides the quantitative model.

% -- NARRATE: What to SAY while showing this slide
Point to the equation: ``Total cost = one-time training + serving cost
times deployment duration times query rate. The serving term scales
with users and time --- it compounds. A 1\% efficiency gain saves
\$1.8M annually for our worked example.'' Walk through the 70B model
example: 1M DAU, 50 req/day, 18B req/yr at \$0.001 = \$18M/yr serving.

% -- ENGAGE: Specific question, prediction, or task for THIS slide
Ask: ``If you can optimize serving by 10\%, how much do you save per
year?'' Expected: 10\% of \$18M = \$1.8M.

% -- WARN: What students will get wrong on THIS topic
Students optimize training efficiency because it is a single system.
Serving efficiency has 10--100x higher leverage because the cost
multiplier is traffic times time.

% -- FLEX: [CORE] or [OPTIONAL] + contingency
[CORE] The cost equation drives all subsequent optimization decisions.
IF SHORT: Show only the equation and the worked example card.
}

\small
$$C_{\text{total}} = \underbrace{C_{\text{training}}}_{\text{one-time}} + \underbrace{C_{\text{serving}} \times T_{\text{deploy}} \times Q_{\text{rate}}}_{\text{ongoing, scales with users}}$$

\vspace{0.1cm}
\begin{mlsyscard}{crimson}
\textbf{Worked Example:} 70B model, 1M DAU, 50 req/day\\[0.05cm]
{\footnotesize
Training: \$2M (one-time) \hfill Serving: 18B req/yr $\times$ \$0.001 = \textbf{\$18M/yr}\\[0.03cm]
Serving costs \textbf{9$\times$} training in year one. A 10\% optimization saves \$1.8M.}
\end{mlsyscard}

\vspace{0.05cm}
\mlsysconcept{Implication:}{Optimize inference ruthlessly, even at the expense of training cost.}

\end{frame}

% =============================================================================
\section{Inversion}
% =============================================================================

\begin{frame}{The Fundamental Inversion: Training vs Inference}
\note{
% -- LINK: What prior concept connects to this slide
The cost equation showed WHY inference matters. This slide shows HOW
inference differs from training at a systems level.

% -- NARRATE: What to SAY while showing this slide
Point to the comparison diagram: ``A training step that takes 2s
instead of 1s is acceptable variance --- throughput averages out. An
inference request that takes 2s instead of 100ms is catastrophic ---
the user sees lag, the SLO is violated, revenue is lost. Training
optimizes throughput. Inference optimizes tail latency.''

% -- ENGAGE: Specific question, prediction, or task for THIS slide
Ask: ``Why does P99 matter more than average latency for inference?''
Expected: because 1 in 100 users experiences the worst case, and at
scale that is thousands of bad experiences per second.

% -- WARN: What students will get wrong on THIS topic
Students measure average latency and report success. P99 is what
users experience during traffic spikes --- the metric that matters for
SLOs and user retention.

% -- FLEX: [CORE] or [OPTIONAL] + contingency
[CORE] The throughput-vs-latency inversion frames every design decision.
IF SHORT: State the inversion and show only the diagram.
}

% --- Layout: FULL-WIDTH diagram ---
\centering
\safeimg[width=0.92\textwidth,height=4.5cm,keepaspectratio]{training-vs-inference.pdf}

\vspace{0.15cm}
\mlsysinsight{Key:}{Training optimizes throughput and tolerates hours of variance. Inference optimizes tail latency and fails at milliseconds.}

\end{frame}

\begin{frame}{The Serving Tax}
\note{
% -- LINK: What prior concept connects to this slide
The inversion showed latency matters. The serving tax quantifies the
overhead that distributed inference adds beyond raw compute.

% -- NARRATE: What to SAY while showing this slide
Walk through the equation: ``Total latency = compute + network +
serialization + coordination + queuing. The non-compute terms consume
10--30\% of the latency budget.'' Point to the table: ``A team
achieving 70 ms on one GPU sees 100 ms in production because they
did not budget for distribution overhead.''

% -- WARN: What students will get wrong on THIS topic
Teams budget only for compute time and are shocked when production
latency is 30\% higher. The serving tax must be accounted for
explicitly in latency budgets.

% -- FLEX: [CORE] or [OPTIONAL] + contingency
[CORE] The serving tax is a key design constraint.
IF SHORT: Show only the equation and the 70 ms to 100 ms example.
}

\small
$$L_{\text{total}} = L_{\text{compute}} + L_{\text{network}} + L_{\text{serial}} + L_{\text{coord}} + L_{\text{queue}}$$

\vspace{0.05cm}
\scriptsize
\renewcommand{\arraystretch}{1.05}
\begin{tabular}{@{}llr@{}}
  \toprule
  \textbf{Component} & \textbf{Source} & \textbf{Typical Range} \\
  \midrule
  \textcolor{computestroke}{$L_{\text{compute}}$} & Model forward pass & 10--500 ms \\
  \textcolor{datastroke}{$L_{\text{network}}$} & Cross-machine RTT & 50--500 $\mu$s \\
  \textcolor{routingstroke}{$L_{\text{serial}} + L_{\text{coord}}$} & Serialization + fan-out & 10--500 $\mu$s \\
  \textcolor{errorstroke}{$L_{\text{queue}}$} & Queuing delay & 0--10 ms \\
  \midrule
  & \textbf{Serving Tax} & \textbf{10--30\% of budget} \\
  \bottomrule
\end{tabular}

\vspace{0.05cm}
\small
\mlsysalert{Budget:}{A team achieving 70 ms on one GPU sees 100 ms in production. Budget distribution overhead explicitly.}

\end{frame}

\begin{frame}{The Serving Hierarchy}
\note{
% -- LINK: What prior concept connects to this slide
The serving tax showed total latency breakdown. The serving hierarchy
organizes optimization by level: request, replica, service, platform.

% -- NARRATE: What to SAY while showing this slide
Point to the diagram then the table: ``Like the memory hierarchy in
computer architecture, each level has different levers. Request level:
batching and caching for per-request latency. Replica level: memory
management and kernel fusion for throughput. Service level: load
balancing and autoscaling for capacity. Platform level: multi-tenancy
for cost efficiency.''

% -- WARN: What students will get wrong on THIS topic
Students optimize at the wrong level. Tuning kernel fusion (replica
level) when the bottleneck is load imbalance (service level) yields
zero improvement.

% -- FLEX: [CORE] or [OPTIONAL] + contingency
[CORE] The hierarchy structures the rest of the chapter.
IF SHORT: Show the hierarchy diagram; skip the detailed table.
}

% --- Layout: FULL-WIDTH diagram ---
\centering
\safeimg[width=0.88\textwidth,height=3.8cm,keepaspectratio]{serving-hierarchy.pdf}

\vspace{0.1cm}
\scriptsize
\renewcommand{\arraystretch}{1.0}
\begin{tabular}{@{}llll@{}}
  \toprule
  \textbf{Level} & \textbf{Target} & \textbf{Key Techniques} & \textbf{Metric} \\
  \midrule
  \textcolor{datastroke}{Request}  & Per-request latency & Batching, caching, prefetch  & ms/request \\
  \textcolor{computestroke}{Replica}  & Throughput     & Memory mgmt, kernel fusion        & tokens/sec \\
  \textcolor{routingstroke}{Service}  & Capacity       & Load balancing, autoscaling       & QPS at SLO \\
  Platform & Efficiency     & Multi-tenancy, scheduling         & \$/query \\
  \bottomrule
\end{tabular}

\end{frame}

% --- ACTIVE LEARNING: Micro-Retrieval Cue ---
\begin{frame}{Quick Check}
\note{
% -- NARRATE: What to SAY while showing this slide
Read question aloud. Wait 15 seconds. Cold-call.
Expected: LLMs are only 1--5\% of production inference by request
volume. RecSys dominate at 80--90\%.

% -- FLEX: [CORE] or [OPTIONAL] + contingency
[CORE] Challenges the LLM-centric assumption before the prediction exercise.
}

\centering
\vspace{1.0cm}
{\Large\bfseries Quick Check}

\vspace{0.6cm}
{\large What fraction of production inference requests\\are LLMs vs.\\ recommendation models?}

\vspace{0.5cm}
{\normalsize\textcolor{midgray}{15 seconds --- then I will cold-call.}}

\end{frame}



% --- ACTIVE LEARNING 1: Predict ---
\begin{frame}{Predict: What Dominates Inference Volume?}
\note{
% -- LINK: What prior concept connects to this slide
The quick check primed the question. This prediction exercise expands
it with options and reveals the answer.

% -- NARRATE: What to SAY while showing this slide
``At major tech companies, which model type accounts for the most
inference requests?'' Give 30 seconds. Most students will guess LLMs
because of media attention.

% -- ENGAGE: Specific question, prediction, or task for THIS slide
Students choose A/B/C/D. After 30 seconds, reveal: C, Recommendation,
80--90\% of requests. Every page load triggers recommendation inference.
LLMs are only 1--5\%.

% -- WARN: What students will get wrong on THIS topic
Students are biased by ChatGPT headlines. LLMs are high-profile but
low-volume. Recommendation systems are invisible but dominate.

% -- FLEX: [CORE] or [OPTIONAL] + contingency
[CORE] Reframes the chapter beyond LLMs.
IF SHORT: Reduce to 15 seconds and reveal immediately.
}

\centering
\vspace{0.4cm}
{\Large\bfseries Think--Write--Share}

\vspace{0.3cm}
{\large At major tech companies (Meta, Google, Netflix),\\
which model type accounts for the \alert{most} inference requests?}

\vspace{0.2cm}
{\normalsize A) LLMs \quad B) Vision (CNN/ViT) \quad C) Recommendation \quad D) Speech}

\vspace{0.2cm}
{\small\textcolor{midgray}{(30 seconds)}}

\pause
\vspace{0.05cm}
\begin{mlsyscard}{datastroke}
{\scriptsize \textbf{C) Recommendation} --- 80--90\% of requests. Every page load triggers inference. LLMs: only 1--5\%.}
\end{mlsyscard}

\end{frame}

% =============================================================================
\section{Batching}
% =============================================================================

\begin{frame}{Why Batching Differs Across Model Types}
\note{
% -- LINK: What prior concept connects to this slide
The prediction exercise showed different model types dominate.
Each type needs a different batching strategy.

% -- NARRATE: What to SAY while showing this slide
Point to the diagram: ``Vision models have fixed input sizes --- static
batching works. LLMs generate variable-length outputs --- continuous
batching is required. RecSys bottleneck is embedding lookup, not dense
compute --- feature-parallel batching applies.''

% -- ENGAGE: Specific question, prediction, or task for THIS slide
Ask: ``Why can we not use the same batch size for ResNet and GPT?''
Expected: ResNet processes fixed-size images in one pass. GPT generates
tokens sequentially with variable output length.

% -- WARN: What students will get wrong on THIS topic
Students assume larger batch = always better. For LLMs, large static
batches waste 60--80\% of compute because all requests wait for the
longest generation.

% -- FLEX: [CORE] or [OPTIONAL] + contingency
[CORE] Establishes the batching taxonomy for the rest of the section.
IF SHORT: Show diagram and state ``no universal optimal.''
}

% --- Layout: FULL-WIDTH diagram ---
\centering
\safeimg[width=0.92\textwidth,height=4.8cm,keepaspectratio]{batching-strategies.pdf}

\vspace{0.1cm}
{\small Match batching strategy to model architecture: \textbf{no universal optimal}.}

\end{frame}

\begin{frame}{The Physics of Batching}
\note{
% -- LINK: What prior concept connects to this slide
The taxonomy showed different batching strategies. This slide provides
the quantitative model.

% -- NARRATE: What to SAY while showing this slide
Walk through the equations: ``Latency is fixed cost plus batch times
variable cost. Throughput is batch over total time. The knee of the
curve is the optimal batch size.'' Show the table: ``Vision models
have low fixed cost (2 ms) so they saturate at B=32. LLMs have high
fixed cost (40 ms, weight loading from HBM) so they need B=128+ to
amortize it.''

% -- WARN: What students will get wrong on THIS topic
Students think the ``optimal'' batch size maximizes throughput. The
real constraint is SLO: find the max B such that L(B) stays within
the latency budget.

% -- FLEX: [CORE] or [OPTIONAL] + contingency
[CORE] The batch size equations are used in exercises and problem sets.
IF SHORT: Show just the equations and the key insight box.
}

\small
\textbf{Latency:} $L(B) = T_{\text{fixed}} + B \times T_{\text{variable}}$

\textbf{Throughput:} $X(B) = \dfrac{B}{T_{\text{fixed}} + B \times T_{\text{variable}}}$

\vspace{0.1cm}
\footnotesize
\renewcommand{\arraystretch}{1.15}
\begin{tabular}{@{}lrrr@{}}
  \toprule
  \textbf{Regime} & \textbf{Vision} & \textbf{LLM} & \textbf{Explanation} \\
  \midrule
  $T_{\text{fixed}}$ & 2 ms & 40 ms & Weight loading from HBM \\
  $T_{\text{variable}}$ & 1.0 ms & 0.5 ms & Per-sample compute \\
  Optimal $B$ (knee) & $\sim$32 & $\sim$128 & Where gains diminish \\
  \bottomrule
\end{tabular}

\vspace{0.1cm}
\mlsysconcept{Engineering goal:}{Find max $B$ such that $L(B) \leq$ SLO. LLMs tolerate larger batches because $T_{\text{fixed}}$ dominates.}

\end{frame}

\begin{frame}{Continuous Batching for LLMs}
\note{
% -- LINK: What prior concept connects to this slide
Static batching works for vision. LLMs need continuous batching
because output lengths vary wildly across requests.

% -- NARRATE: What to SAY while showing this slide
Walk through the problem: ``Static batching pads all requests to the
longest generation. If one request generates 10 tokens and another
generates 500, the short request wastes 490 token slots --- 60--80\%
wasted compute.'' Then the solution: ``Continuous batching lets
requests enter and exit at each decode step. When one request
finishes, a new request fills the slot. vLLM achieved 2--4x throughput
over FasterTransformer.''

% -- ENGAGE: Specific question, prediction, or task for THIS slide
Ask: ``Does continuous batching eliminate ALL waste?'' Expected: no ---
prefill is still compute-bound (O(n\^2) attention) and requests
cannot share that compute. Continuous batching addresses decode waste
only.

% -- WARN: What students will get wrong on THIS topic
Students think continuous batching solves all LLM serving problems.
It addresses decode throughput but not prefill latency, KV cache
limits, or sharding bandwidth.

% -- FLEX: [CORE] or [OPTIONAL] + contingency
[CORE] Continuous batching is the most important serving optimization for LLMs.
IF SHORT: Focus on the 60--80\% waste problem and the 2--4x result.
}

\small
\begin{columns}[T]
  \begin{column}{0.50\textwidth}
    \textbf{Problem:} Static batching for LLMs

    \vspace{0.1cm}
    {\footnotesize
    \begin{itemize}\setlength\itemsep{0pt}
      \item Each request generates different \# tokens
      \item All requests wait for the \emph{longest} generation
      \item GPU idle time: \textbf{60--80\%} wasted compute
    \end{itemize}
    }

    \vspace{0.15cm}
    \textbf{Solution:} Continuous batching

    \vspace{0.1cm}
    {\footnotesize
    \begin{itemize}\setlength\itemsep{0pt}
      \item Requests enter/exit at each decode step
      \item Finished request $\to$ new request fills slot
      \item \textbf{2--4$\times$} throughput over static batching
    \end{itemize}
    }
  \end{column}
  \begin{column}{0.46\textwidth}
    \begin{mlsyscard}{datastroke}
    \textbf{vLLM (2023):}\\
    {\footnotesize OS-style virtual memory for KV cache.
    2--4$\times$ over FasterTransformer,
    up to 24$\times$ over HuggingFace,
    by eliminating 60--80\% memory fragmentation.}
    \end{mlsyscard}

    \vspace{0.15cm}
    \begin{mlsyscard}{routingstroke}
    \textbf{Prefill vs.\ Decode:}\\
    {\footnotesize Prefill = compute-bound ($O(n^2)$).\\
    Decode = memory-bandwidth-bound.\\
    Two distinct optimization regimes.}
    \end{mlsyscard}
  \end{column}
\end{columns}

\end{frame}

\begin{frame}{Queuing Theory for Batched Inference}
\note{
% -- LINK: What prior concept connects to this slide
Continuous batching increased throughput. But how do we model the interaction between arrival rate, batch size, and latency? Queuing theory provides the framework.

% -- NARRATE: What to SAY while showing this slide
``LLM serving is an M/G/c/K queue: Poisson arrivals (M), general service time distribution (G) because generation lengths vary, c GPU replicas, and K maximum queue depth. The key result: utilization above 70\% causes P99 latency to explode exponentially. For GPT-3 at 100 QPS on 8 H100s: mean service time is 200 ms with batch=8, utilization is 100*0.2/8 = 0.25 --- comfortable. But at 400 QPS: utilization hits 1.0 and the queue grows unboundedly.''

% -- ENGAGE: Specific question, prediction, or task for THIS slide
Ask: ``At what QPS does P99 latency double for this setup?'' Expected: around 280--300 QPS (utilization ~0.7), where the queuing delay hockey stick begins.

% -- WARN: What students will get wrong on THIS topic
Students think latency scales linearly with QPS. It scales exponentially above 70\% utilization. This is why production systems target 50--60\% GPU utilization.

% -- FLEX: [OPTIONAL] Mathematical foundation for capacity planning.
IF SHORT: State the 70\% utilization threshold and skip the queuing model details.
}

\footnotesize
\begin{columns}[T]
  \begin{column}{0.50\textwidth}
    \textbf{M/G/c/K Queuing Model:}

    \vspace{0.1cm}
    {\scriptsize
    \begin{tabular}{@{}ll@{}}
      $\lambda$ & Arrival rate (QPS) \\
      $\mu$ & Service rate per GPU ($1/T_{\text{gen}}$) \\
      $c$ & Number of GPU replicas \\
      $\rho = \lambda / (c \cdot \mu)$ & Utilization \\
    \end{tabular}
    }

    \vspace{0.15cm}
    \begin{mlsyscard}{errorstroke}
    {\scriptsize At $\rho > 0.7$: P99 latency explodes.\\
    Production target: \textbf{50--60\% utilization}.}
    \end{mlsyscard}
  \end{column}
  \begin{column}{0.46\textwidth}
    \textbf{GPT-3 at 100 QPS, 8$\times$ H100:}

    \vspace{0.1cm}
    {\scriptsize
    \renewcommand{\arraystretch}{1.1}
    \begin{tabular}{@{}lr@{}}
      \toprule
      Mean service time & 200 ms \\
      Service rate $\mu$ & 5 req/s/GPU \\
      Utilization $\rho$ & $100/(8 \times 5) = 0.25$ \\
      P99 latency & $\sim$400 ms \\
      \midrule
      \textbf{At 400 QPS:} & \\
      Utilization $\rho$ & $\to 1.0$ \\
      P99 latency & $\to \infty$ \\
      \bottomrule
    \end{tabular}
    }
  \end{column}
\end{columns}

\end{frame}

% =============================================================================
\section{KV Cache}
% =============================================================================

\mlsysfocus{The KV Cache Wall}{%
$\text{KV size} = 2 \times L \times H \times S \times B \times P$\\[0.3cm]
{\normalsize layers $\times$ heads $\times$ seq\_len $\times$ batch $\times$ precision}%
}

\begin{frame}{KV Cache: Memory Exhaustion in LLM Serving}
\note{
% -- LINK: What prior concept connects to this slide
Continuous batching increased throughput but the KV cache limits how
many requests can run concurrently. This is the memory wall of serving.

% -- NARRATE: What to SAY while showing this slide
Point to the figure: ``Model weights are fixed at 35 GB for a 70B
model in INT4. But the KV cache grows linearly with batch size times
sequence length. At batch 8 with 48K tokens, you hit 80 GB --- the
full H100.'' Walk through the bottom columns: ``Weights are fixed.
KV cache is the variable that exhausts memory.''

% -- ENGAGE: Specific question, prediction, or task for THIS slide
Ask: ``Why not just add more memory?'' Expected: the bottleneck shifts
to bandwidth before memory is exhausted. Reading a large KV cache
from HBM becomes the latency bottleneck.

% -- WARN: What students will get wrong on THIS topic
Students think KV cache is small because each token's KV is only a
few hundred KB. But multiplied by layers, heads, batch, and sequence
length, it quickly exceeds model weights.

% -- FLEX: [CORE] or [OPTIONAL] + contingency
[CORE] The KV cache wall motivates PagedAttention on the next slide.
IF SHORT: Show the figure and state the B=8 at 48K = OOM result.
}

% --- Layout: FULL-WIDTH diagram ---
\centering
\safeimg[width=0.88\textwidth,height=4.5cm,keepaspectratio]{kv-cache-wall.pdf}

\vspace{0.1cm}
\small
\begin{columns}[T]
  \begin{column}{0.45\textwidth}
    \centering
    \textcolor{computestroke}{\textbf{Model Weights}}: Fixed 35 GB\\
    {\scriptsize (70B params, INT4 quantized)}
  \end{column}
  \begin{column}{0.45\textwidth}
    \centering
    \textcolor{errorstroke}{\textbf{KV Cache}}: Grows with $B \times S$\\
    {\scriptsize B=8 hits OOM at 48K tokens}
  \end{column}
\end{columns}

\end{frame}

\begin{frame}{PagedAttention: Virtual Memory for KV Cache}
\note{
% -- LINK: What prior concept connects to this slide
The KV cache wall showed memory exhaustion. PagedAttention is the
solution, borrowed directly from OS virtual memory.

% -- NARRATE: What to SAY while showing this slide
``Standard serving pre-allocates max sequence length per request ---
if max is 4096 tokens but the request uses 200, 95\% of the
allocation is wasted. PagedAttention allocates pages on demand, just
like virtual memory in an OS. Pages are 16--64 tokens. This recovers
60--80\% of wasted memory, enabling 2--4x more concurrent requests on
the same hardware.''

% -- ENGAGE: Specific question, prediction, or task for THIS slide
Ask: ``What OS concept does this remind you of?'' Expected: virtual
memory with demand paging. vLLM's insight was recognizing that the
KV cache allocation problem is isomorphic to OS page management.

% -- WARN: What students will get wrong on THIS topic
Students think PagedAttention changes the attention algorithm. It does
not --- the math is identical. It only changes how KV cache memory is
allocated, using a page table indirection.

% -- FLEX: [CORE] or [OPTIONAL] + contingency
[CORE] PagedAttention is the defining innovation of modern LLM serving.
IF SHORT: State the OS analogy and the 2--4x result.
}

% --- Layout: FULL-WIDTH diagram ---
\centering
\safeimg[width=0.92\textwidth,height=5.0cm,keepaspectratio]{paged-attention.pdf}

\vspace{0.1cm}
\mlsysinsight{Result:}{Near-100\% memory utilization. 2--4$\times$ more concurrent requests on the same hardware.}

\end{frame}

% --- ACTIVE LEARNING 2: Exercise ---
\begin{frame}{Your Turn: KV Cache Memory Budget}
\note{
% -- LINK: What prior concept connects to this slide
The KV cache wall showed the problem. This exercise applies the
formula to a real model and checks if it fits on hardware.

% -- NARRATE: What to SAY while showing this slide
``Llama-70B: 80 layers, 8 KV heads (GQA), dim=128, FP16, batch=16,
seq=4096. Calculate the KV cache size and check if it fits on an 80 GB
H100.'' Give 90 seconds.

% -- ENGAGE: Specific question, prediction, or task for THIS slide
Solution: 2 * 80 * 8 * 128 * 2 = 327K bytes/token. Times 4096 * 16
= 21.5 GB. Weights 35 + KV 21.5 = 56.5 GB --- fits. But at B=32:
KV = 43 GB, total = 78 GB, barely fits.

% -- WARN: What students will get wrong on THIS topic
Students forget the factor of 2 (K and V are separate). They also
confuse KV heads (8 for GQA) with attention heads (64). GQA reduces
KV cache by 8x.

% -- FLEX: [CORE] or [OPTIONAL] + contingency
[CORE] Quantitative exercise testing the KV cache formula.
IF SHORT: Reduce to 45 seconds.
}

\footnotesize
\begin{columns}[T]
  \begin{column}{0.55\textwidth}
    \textbf{Calculate KV Cache Memory}

    \vspace{0.05cm}
    {\scriptsize Llama-70B: $L\!=\!80$, $H\!=\!8$ KV heads, dim$\!=\!128$, FP16. $B\!=\!16$, $S\!=\!4096$.}

    \vspace{0.05cm}
    KV cache size (GB)? Fits on 80 GB H100? {\scriptsize\textcolor{midgray}{(90 sec)}}
  \end{column}
  \begin{column}{0.42\textwidth}
    \pause
    \begin{mlsyscard}{computestroke}
    {\scriptsize \textbf{Solution:}\\
    $2 \!\times\! 80 \!\times\! 8 \!\times\! 128 \!\times\! 2$ = 327K B/tok\\
    $\times 4096 \times 16 =$ \textbf{21.5 GB}\\
    Weights 35 + KV 21.5 = \textbf{56.5 GB} --- fits!\\
    \textcolor{errorstroke}{B=32: KV=43 $\to$ 78 GB (barely!)}}
    \end{mlsyscard}
  \end{column}
\end{columns}

\end{frame}

\begin{frame}{Speculative Decoding}
\note{
% -- LINK: What prior concept connects to this slide
Continuous batching and PagedAttention optimize throughput. Speculative
decoding optimizes latency by generating multiple tokens per step.

% -- NARRATE: What to SAY while showing this slide
``LLM decode is sequential --- each token waits for the previous. The
insight: use a small draft model (1--7B) to generate k candidate tokens
fast, then verify all k in one pass with the large model. Accepted
tokens are mathematically identical to sequential generation. Rejected
tokens waste compute, but compute is abundant (memory is the bottleneck).
Result: 2--3x speedup for greedy generation.''

% -- WARN: What students will get wrong on THIS topic
Students think speculative decoding is approximate. It is exact ---
the verification step ensures output is identical to sequential
generation (for greedy decoding).

% -- FLEX: [OPTIONAL] + contingency
[OPTIONAL] Speculative decoding is covered more deeply in Ch 10.
IF SHORT: State the insight and 2--3x result; skip the requirements box.
}

\footnotesize
\begin{columns}[T]
  \begin{column}{0.55\textwidth}
    \textbf{Problem:} LLM decode is sequential --- each token waits for the previous.

    \vspace{0.1cm}
    \textbf{Insight:} Draft $\to$ Verify in parallel

    \vspace{0.05cm}
    {\footnotesize
    \begin{enumerate}\setlength\itemsep{0pt}
      \item Small \textbf{draft model} (1--7B) generates $k$ candidates
      \item Large \textbf{target model} (70B+) verifies all $k$ at once
      \item Accept matches, regenerate from first divergence
    \end{enumerate}
    }

    \vspace{0.1cm}
    \mlsysconcept{Speedup:}{2--3$\times$ for greedy generation when acceptance rate $> 70\%$.}
  \end{column}
  \begin{column}{0.42\textwidth}
    \begin{mlsyscard}{datastroke}
    {\footnotesize \textbf{Requirements:} Same vocabulary. Draft model must be fast enough that draft + verify $<$ sequential decode.}
    \end{mlsyscard}

    \vspace{0.1cm}
    \begin{mlsyscard}{errorstroke}
    {\footnotesize \textbf{Trade-off:} Extra memory for draft model. Wasted compute on rejected tokens. Best for deterministic outputs.}
    \end{mlsyscard}
  \end{column}
\end{columns}

\end{frame}

% =============================================================================
\section{Sharding}
% =============================================================================

\begin{frame}{Model Sharding for Inference}
\note{
% -- LINK: What prior concept connects to this slide
KV cache and batching optimize single-GPU serving. When models exceed
GPU memory, sharding across multiple GPUs is mandatory.

% -- NARRATE: What to SAY while showing this slide
Walk through the table: ``Llama-70B in FP16 = 140 GB, exceeding the
80 GB H100. Minimum 2-way sharding. For long-context models, KV cache
alone can force 4--8 way sharding. For recommendation systems with
100 TB embedding tables, 1000+ way sharding.'' Point to the equation:
``Parallel time = sequential time / P + communication. Sharding helps
only when communication overhead is less than time saved.''

% -- ENGAGE: Specific question, prediction, or task for THIS slide
Ask: ``What parallelism strategies from distributed training apply
here?'' Expected: tensor and pipeline parallelism. The key difference:
latency matters for inference, not throughput.

% -- WARN: What students will get wrong on THIS topic
Students apply training sharding strategies directly to inference.
Pipeline parallelism creates pipeline bubbles that are acceptable for
training throughput but unacceptable for inference latency. Tensor
parallelism is preferred for latency-critical serving.

% -- FLEX: [CORE] or [OPTIONAL] + contingency
[CORE] Sharding is required for all frontier models.
IF SHORT: Show only the table and the constraint equation.
}

\small
\textbf{When sharding becomes necessary:}

\vspace{0.1cm}
\footnotesize
\renewcommand{\arraystretch}{1.15}
\begin{tabular}{@{}llrl@{}}
  \toprule
  \textbf{Trigger} & \textbf{Model Example} & \textbf{Min Sharding} & \textbf{Strategy} \\
  \midrule
  Memory (weights) & Llama-70B (140 GB) & 2-way & Tensor/pipeline \\
  Memory (KV cache) & GPT-4 (long ctx) & 4--8 way & Tensor (cache) \\
  Memory (embeddings) & DLRM (100 TB) & 1000+ way & Embedding shard \\
  Latency & Any large model & Varies & Tensor parallel \\
  \bottomrule
\end{tabular}

\vspace{0.15cm}
\small
$$T_{\text{parallel}} = \frac{T_{\text{sequential}}}{P} + T_{\text{communication}}$$

\vspace{0.1cm}
\mlsysalert{Constraint:}{Sharding helps only when communication overhead $<$ time saved. Requires NVLink (900 GB/s), not PCIe (64 GB/s).}

\end{frame}

\begin{frame}{Tensor Parallelism vs Pipeline Parallelism}
\note{
% -- LINK: What prior concept connects to this slide
Sharding is necessary. The question is which strategy: tensor or
pipeline parallelism.

% -- NARRATE: What to SAY while showing this slide
Left card: ``Tensor parallelism splits each layer across GPUs.
Attention heads are partitioned: H/P per device. Two AllReduces per
transformer block. Lowest latency but highest bandwidth requirement
--- must use NVLink.'' Right card: ``Pipeline parallelism splits the
model into P sequential stages. Lower bandwidth needs but pipeline
bubbles increase latency.''

% -- ENGAGE: Specific question, prediction, or task for THIS slide
Ask: ``For latency-critical serving, which would you choose?''
Expected: tensor parallelism --- it has no pipeline bubbles.

% -- WARN: What students will get wrong on THIS topic
Students use pipeline parallelism for inference because it works
across machines (no NVLink needed). But the pipeline bubbles create
unacceptable latency variance for serving SLOs.

% -- FLEX: [CORE] or [OPTIONAL] + contingency
[CORE] The TP-vs-PP decision determines serving architecture.
IF SHORT: State the rule --- TP within node, PP across nodes.
}

\scriptsize
\begin{columns}[T]
  \begin{column}{0.48\textwidth}
    \begin{mlsyscard}{computestroke}
    {\scriptsize \textbf{Tensor Parallelism}\\
    Split each layer across GPUs. Heads partitioned: $H/P$ per device. AllReduce per block.}
    \end{mlsyscard}

    \vspace{0.05cm}
    \textbf{Pros:} Lowest latency, no pipeline bubbles\\
    \textbf{Cons:} 2 AllReduce/layer, needs NVLink
  \end{column}
  \begin{column}{0.48\textwidth}
    \begin{mlsyscard}{routingstroke}
    {\scriptsize \textbf{Pipeline Parallelism}\\
    Split model into $P$ stages. Each GPU processes a layer group. Point-to-point between stages.}
    \end{mlsyscard}

    \vspace{0.05cm}
    \textbf{Pros:} Lower BW needs, works cross-machine\\
    \textbf{Cons:} Pipeline bubbles, higher latency
  \end{column}
\end{columns}

\vspace{0.1cm}
\small
\mlsysconcept{Rule:}{Tensor parallelism within a node (NVLink). Pipeline parallelism across nodes (network).}

\end{frame}

\begin{frame}{Serving Architecture Dimensions}
\note{
% -- LINK: What prior concept connects to this slide
Sharding and batching are individual techniques. This slide presents the full decision space for serving architecture design as a structured framework.

% -- NARRATE: What to SAY while showing this slide
``Four dimensions define your serving architecture. Scheduling policy: FCFS for fairness, preemptive for SLO compliance. Deployment topology: single-GPU for small models, tensor-parallel for latency, pipeline for throughput. State management: stateless for simple scaling, stateful for KV cache reuse across turns. Compute allocation: dedicated GPUs for predictable latency, shared for cost efficiency.''

% -- ENGAGE: Specific question, prediction, or task for THIS slide
Ask: ``For a chatbot with multi-turn conversations, which state management model would you choose?'' Expected: stateful, because KV cache from previous turns can be reused, avoiding redundant prefill.

% -- WARN: What students will get wrong on THIS topic
Students default to stateless because it is simpler to scale. But stateful serving with KV cache persistence can reduce prefill cost by 60--80\% for multi-turn conversations.

% -- FLEX: [OPTIONAL] Framework for serving architecture decisions.
IF SHORT: Point to the table and state the key trade-off.
}

\footnotesize
\renewcommand{\arraystretch}{1.15}
\begin{tabular}{@{}lp{3.5cm}p{4.5cm}@{}}
  \toprule
  \textbf{Dimension} & \textbf{Options} & \textbf{Trade-off} \\
  \midrule
  \textbf{Scheduling} & FCFS, Preemptive, Priority & Fairness vs.\ SLO compliance \\
  \textbf{Topology}   & Single-GPU, TP, PP, TP+PP & Latency vs.\ throughput vs.\ cost \\
  \textbf{State}      & Stateless, Stateful (KV reuse) & Scaling simplicity vs.\ cache efficiency \\
  \textbf{Compute}    & Dedicated, Shared, Spot & Predictability vs.\ cost \\
  \bottomrule
\end{tabular}

\vspace{0.15cm}
\begin{mlsyscard}{routingstroke}
{\scriptsize\textbf{Multi-turn chatbot:} Stateful + TP + Preemptive. KV cache reuse saves 60--80\% prefill cost across conversation turns.}
\end{mlsyscard}

\end{frame}

% =============================================================================
\section{Load Balancing}
% =============================================================================

\begin{frame}{Load Balancing for Inference}
\note{
% -- LINK: What prior concept connects to this slide
Sharding placed the model. Load balancing distributes requests across
replicas of the model.

% -- NARRATE: What to SAY while showing this slide
``Simple round-robin fails for generative ML because request processing
times vary by 10x or more --- a 10-token response vs. a 1000-token
response. The power-of-two-choices result is one of the most impactful
findings in distributed systems: probe two random servers instead of
one, and pick the less loaded. This provides exponential tail latency
reduction with only one extra probe.''

% -- ENGAGE: Specific question, prediction, or task for THIS slide
Ask: ``Why does round-robin fail for LLM serving?'' Expected: because
generation lengths vary 10--100x, so a ``balanced'' assignment creates
massive imbalance in actual load.

% -- WARN: What students will get wrong on THIS topic
Students think round-robin is ``fair enough.'' For uniform-cost
requests it is. For variable-cost LLM requests, round-robin creates
hotspots that violate P99 SLOs.

% -- FLEX: [CORE] or [OPTIONAL] + contingency
[CORE] Power-of-two-choices is a foundational distributed systems result.
IF SHORT: Show diagram and state the exponential improvement.
}

% --- Layout: FULL-WIDTH diagram ---
\centering
\safeimg[width=0.92\textwidth,height=4.8cm,keepaspectratio]{power-of-two-choices.pdf}

\vspace{0.1cm}
{\small One extra probe per request provides \textbf{exponential} tail latency reduction.}

\end{frame}

\begin{frame}{Power of Two Choices: The Math}
\note{
% -- LINK: What prior concept connects to this slide
The previous slide showed the result visually. This slide provides
the mathematical proof of the exponential improvement.

% -- NARRATE: What to SAY while showing this slide
``Random assignment: max queue grows as log n / log log n. Two choices:
max queue grows as log log n. For 1000 servers: random has max queue
of about 5, two choices has about 2. The key: with random routing,
unlucky picks compound. With two choices, you avoid the worst pick
every time --- the probability of landing on a bad server drops from
d/n to (d/n)\^2.''

% -- WARN: What students will get wrong on THIS topic
Students think you need sophisticated algorithms (least-loaded,
weighted round-robin). Two random choices achieves most of the
theoretical benefit with zero state tracking.

% -- FLEX: [OPTIONAL] + contingency
[OPTIONAL] Mathematical detail supplements the intuitive diagram.
IF SHORT: Skip; the previous slide is sufficient.
}

\footnotesize
\textbf{Random assignment:} \quad
$E[\text{max queue}] = \Theta\!\left(\frac{\log n}{\log \log n}\right)$

\vspace{0.05cm}
\textbf{Power of two choices:} \quad
$E[\text{max queue}] = \Theta(\log \log n)$

\vspace{0.1cm}
\scriptsize
\renewcommand{\arraystretch}{1.05}
\begin{tabular}{@{}lrr@{}}
  \toprule
  \textbf{Servers ($n$)} & \textbf{Random} & \textbf{Two Choices} \\
  \midrule
  100   & $\sim$3--4  & $\sim$2 \\
  1,000 & $\sim$4--5  & $\sim$2 \\
  10,000 & $\sim$5--6 & $\sim$2--3 \\
  \bottomrule
\end{tabular}

\vspace{0.05cm}
\footnotesize
\begin{mlsyscard}{datastroke}
{\scriptsize \textbf{Why:} Random routing occasionally picks busy servers; mistakes compound. Two choices avoids worst picks --- $(d/n)^2$ decay.}
\end{mlsyscard}
\mlsyscite{Mitzenmacher, ``Power of Two Choices,'' 2001}

\end{frame}

% --- ACTIVE LEARNING: Power of Two Choices Exercise ---
\begin{frame}{Your Turn: Power of Two Choices at Scale}
\note{
% -- LINK: What prior concept connects to this slide
Students just saw the P2C math. This exercise makes them compute the concrete improvement for a realistic serving fleet.

% -- NARRATE: What to SAY while showing this slide
Read the problem. Give 90 seconds. Solution: ``Random: expected max load = log(1000)/log(log(1000)) = 6.9/1.93 = 3.6 requests above average. P2C: expected max load = log(log(1000)) = log(6.9) = 1.93 above average. Improvement: 3.6/1.93 = 1.9x reduction in max load. For P99 with 50 ms mean: random P99 = 50 * 3.6 = 180 ms. P2C P99 = 50 * 1.93 = 97 ms. Nearly 2x improvement from one extra probe.''

% -- ENGAGE: Specific question, prediction, or task for THIS slide
After revealing, ask: ``Would power of three choices be worth the extra probe?'' Expected: marginal --- P2C captures most of the benefit. Each additional choice gives diminishing returns.

% -- WARN: What students will get wrong on THIS topic
Students will try to compute exact P99 values. The log/log-log expressions give the scaling behavior, not exact percentiles. Emphasize the ratio, not the absolute numbers.

% -- FLEX: [CORE] Quantitative exercise applying P2C to a real serving scenario.
IF SHORT: Show the solution and state the 1.9x improvement.
}

\small
\begin{columns}[T]
  \begin{column}{0.55\textwidth}
    {\normalsize\bfseries Compute P2C Improvement}

    \vspace{0.15cm}
    \textbf{Setup:} 1,000 requests across 100 servers.\\
    Mean service time: 50 ms.

    \vspace{0.1cm}
    \begin{enumerate}
      \item Expected max load with random routing?
      \item Expected max load with P2C?
      \item P99 latency improvement factor?
    \end{enumerate}

    {\footnotesize\textcolor{midgray}{(90 seconds --- use the formulas from the previous slide)}}
  \end{column}
  \begin{column}{0.42\textwidth}
    \pause
    \begin{mlsyscard}{datastroke}
    \textbf{Solution:}\\[0.1cm]
    {\footnotesize
    Random: $\frac{\log 1000}{\log\log 1000} \approx$ \textbf{3.6}\\
    P2C: $\log\log 1000 \approx$ \textbf{1.9}\\[0.05cm]
    P99 ratio: \textbf{1.9$\times$ improvement}\\[0.1cm]
    Random P99: $\sim$180 ms\\
    P2C P99: $\sim$97 ms\\[0.05cm]
    \textcolor{datastroke}{\textbf{One extra probe = 2$\times$ tail latency reduction.}}
    }
    \end{mlsyscard}
  \end{column}
\end{columns}

\end{frame}

% --- ACTIVE LEARNING 3: Discussion ---
\begin{frame}{Discussion: Which Level Is the Bottleneck?}
\note{
% -- LINK: What prior concept connects to this slide
Students have seen the serving hierarchy (request, replica, service,
platform). This discussion applies it to diagnose a real scenario.

% -- NARRATE: What to SAY while showing this slide
``An LLM API serves 10K QPS with P99 = 2 seconds. Where is the
bottleneck?'' Turn-and-talk 90 seconds. Cold-call 2--3 pairs.

% -- ENGAGE: Specific question, prediction, or task for THIS slide
Students classify the bottleneck by hierarchy level. No single right
answer --- it depends on the workload. LLMs are often request-level
(KV cache). RecSys are service-level (load imbalance).

% -- WARN: What students will get wrong on THIS topic
Students jump to ``add more GPUs'' (platform level) without diagnosing
whether the bottleneck is batching (request), KV cache (replica), or
load imbalance (service). Diagnosis before prescription.

% -- FLEX: [CORE] or [OPTIONAL] + contingency
[CORE] Synthesizes the serving hierarchy into a diagnostic exercise.
IF SHORT: Do show-of-hands for each level instead of turn-and-talk.
}

\centering
\vspace{0.6cm}
{\large\bfseries Turn and Talk \textcolor{midgray}{(90 seconds)}}

\vspace{0.5cm}
{\large An LLM API serves 10K QPS with P99 = 2 seconds.\\[0.2cm]
\alert{At which serving hierarchy level is the bottleneck?}}

\vspace{0.4cm}
\begin{columns}[c]
  \begin{column}{0.22\textwidth}\centering\small\textcolor{datastroke}{\textbf{Request}}\\Batching?\end{column}
  \begin{column}{0.22\textwidth}\centering\small\textcolor{computestroke}{\textbf{Replica}}\\KV Cache?\end{column}
  \begin{column}{0.22\textwidth}\centering\small\textcolor{routingstroke}{\textbf{Service}}\\Load Bal.?\end{column}
  \begin{column}{0.22\textwidth}\centering\small\textbf{Platform}\\Resources?\end{column}
\end{columns}

\end{frame}

% =============================================================================
\section{Autoscaling}
% =============================================================================

\begin{frame}{The Cold Start Problem}
\note{
% -- LINK: What prior concept connects to this slide
Load balancing distributes traffic across existing replicas. Autoscaling
adds new replicas. But GPU cold start is 100x slower than web services.

% -- NARRATE: What to SAY while showing this slide
Point to the timeline: ``A web service cold start takes seconds. A 70B
model takes 5+ minutes to provision a VM, load weights, and warm
caches. Reactive scaling alone cannot handle traffic spikes because
the new capacity arrives minutes after the SLO violation.''

% -- WARN: What students will get wrong on THIS topic
Students assume autoscaling works as fast as for web services. GPU
cold start (provisioning + weight loading + warmup) makes reactive
scaling insufficient. Predictive scaling and warm pools are mandatory.

% -- FLEX: [CORE] or [OPTIONAL] + contingency
[CORE] Cold start is the fundamental constraint of inference autoscaling.
IF SHORT: Show timeline and state ``5+ min cold start means reactive
scaling fails.''
}

% --- Layout: FULL-WIDTH diagram ---
\centering
\safeimg[width=0.92\textwidth,height=3.5cm,keepaspectratio]{cold-start-timeline.pdf}

\vspace{0.05cm}
\scriptsize
\renewcommand{\arraystretch}{0.9}
\begin{tabular}{@{}lrl@{}}
  \toprule
  \textbf{Scaling Type} & \textbf{Response} & \textbf{Trade-off} \\
  \midrule
  Batch size & Instant & Increases per-request latency \\
  Horizontal & Minutes & Linear cost, slow cold start \\
  Vertical & Hours & Non-linear cost, redeployment \\
  \bottomrule
\end{tabular}

\vspace{0.05cm}
\small
\mlsysalert{Key:}{Scaling decisions must anticipate demand 5+ min ahead. Batch size is the first knob to turn.}

\end{frame}

\begin{frame}{Autoscaling Strategies}
\note{
% -- LINK: What prior concept connects to this slide
Cold start showed the constraint. This slide presents the two
complementary strategies: reactive and predictive.

% -- NARRATE: What to SAY while showing this slide
Left side: ``Reactive scaling monitors queue depth, GPU util, P99.
Scale up when metrics cross thresholds. Risk: 5 min cold start
exceeds SLO violation window.'' Right side: ``Predictive scaling
forecasts traffic 5--10 min ahead. Warm pools keep pre-loaded replicas
ready to serve.'' Point to the cold start card: ``7B: 30--60s. 70B:
5+ min. 175B: 12+ min. The larger the model, the more critical
predictive scaling becomes.''

% -- WARN: What students will get wrong on THIS topic
Students implement only reactive scaling because it is simpler.
Without predictive scaling or warm pools, every traffic spike causes
SLO violations during the 5+ minute cold start window.

% -- FLEX: [CORE] or [OPTIONAL] + contingency
[CORE] Autoscaling completes the serving system design.
IF SHORT: Focus on the cold start times and the ``combine both'' conclusion.
}

\footnotesize
\begin{columns}[T]
  \begin{column}{0.48\textwidth}
    \begin{mlsyscard}{routingstroke}
    {\scriptsize \textbf{Reactive Scaling}\\
    Monitor: GPU util, queue depth, P99.
    Scale up when queue $>$ threshold.\\
    \textcolor{errorstroke}{Risk: 5 min cold start $>$ SLO violation}}
    \end{mlsyscard}

    \vspace{0.1cm}
    \begin{mlsyscard}{datastroke}
    {\scriptsize \textbf{Predictive Scaling}\\
    Forecast traffic 5--10 min ahead.
    \textcolor{datastroke}{Warm pools}: pre-loaded replicas ready to serve.}
    \end{mlsyscard}
  \end{column}
  \begin{column}{0.48\textwidth}
    \begin{mlsyscard}{crimson}
    {\scriptsize \textbf{Cold Start:} $T_{\text{cold}} = T_{\text{prov}} + T_{\text{load}} + T_{\text{warm}}$\\[0.05cm]
    7B: $\sim$30--60s \quad 70B: $\sim$5+ min \quad 175B: $\sim$12+ min}
    \end{mlsyscard}

    \vspace{0.1cm}
    \mlsysconcept{Production:}{Combine reactive + predictive. Warm pools absorb spikes. Batch size buys time.}
  \end{column}
\end{columns}

\end{frame}

% =============================================================================
\section{Wrap-Up}
% =============================================================================

\begin{frame}{Fallacies}
\note{
% -- LINK: What prior concept connects to this slide
The chapter covered economics, batching, KV cache, sharding, load
balancing, and autoscaling. These fallacies challenge remaining
misconceptions.

% -- NARRATE: What to SAY while showing this slide
Read each fallacy and rebuttal: ``LLMs are only 1--5\% of inference
volume. Continuous batching solves decode throughput but not prefill,
KV cache, or sharding. Power-of-two-choices gives exponential
improvement with one extra probe.''

% -- FLEX: [CORE] or [OPTIONAL] + contingency
[CORE] Quantitative rebuttals reinforce key numbers.
IF SHORT: Cover only the first fallacy (LLMs vs. RecSys).
}

\footnotesize
\textbf{Fallacy:} \textit{Inference at scale is synonymous with LLM serving.}\\
{\scriptsize RecSys are 80--90\% of production inference by request volume. LLMs are only 1--5\%. Feature-parallel batching and embedding sharding dominate production, not continuous batching.}

\vspace{0.08cm}
\textbf{Fallacy:} \textit{Continuous batching solves all LLM serving problems.}\\
{\scriptsize Addresses decode throughput but not prefill ($O(n^2)$ attention), KV cache memory limits, or network bandwidth for sharded models. Necessary but not sufficient.}

\vspace{0.08cm}
\textbf{Fallacy:} \textit{Load balancing does not matter for inference.}\\
{\scriptsize Power-of-two-choices: max queue from $O(\log n / \log \log n)$ to $O(\log \log n)$. For 1,000 servers: $\sim$5 to $\sim$2. First-order effect on P99.}

\end{frame}

\begin{frame}{Pitfalls}
\note{
% -- LINK: What prior concept connects to this slide
Fallacies addressed beliefs. Pitfalls address operational mistakes.

% -- NARRATE: What to SAY while showing this slide
``Using training infrastructure for serving discovers unacceptable
latency variance. Sizing on average throughput means a 25\% spike
causes unbounded queues. Ignoring the serving tax means production
latency is 30\% higher than benchmarks.''

% -- FLEX: [CORE] or [OPTIONAL] + contingency
[CORE] Operational pitfalls for serving systems.
IF SHORT: Cover just the first two pitfalls.
}

\footnotesize
\textbf{Pitfall:} \textit{Using training infrastructure for production serving.}\\
{\scriptsize Training optimizes throughput over hours; serving optimizes P99 in milliseconds. Teams deploying training clusters for serving discover unacceptable latency variance.}

\vspace{0.08cm}
\textbf{Pitfall:} \textit{Sizing capacity based on average throughput.}\\
{\scriptsize At 80\% utilization, a 25\% spike causes unbounded queue growth. GPU cold start (5+ min) means capacity arrives after SLO violations. Plan for peak + headroom.}

\vspace{0.08cm}
\textbf{Pitfall:} \textit{Ignoring the serving tax in latency budgets.}\\
{\scriptsize Distributed inference adds 10--30\% overhead. A team achieving 70 ms on one GPU sees 100 ms in production.}

\end{frame}

% --- RETRIEVAL PRACTICE ---

% --- MUDDIEST POINT ---
\begin{frame}{Muddiest Point}
\note{
% -- NARRATE: What to SAY while showing this slide
``Write down the most confusing concept from today. Anonymous, one
sentence. I will address top confusions next lecture.''

% -- FLEX: [CORE] or [OPTIONAL] + contingency
[CORE] Feedback loop.
}

\centering
\vspace{1.0cm}
{\Large\bfseries What was the \alert{muddiest point} today?}

\vspace{0.8cm}
{\normalsize Write down the concept you found \textbf{most confusing}.}

\vspace{0.5cm}
{\small\textcolor{midgray}{Anonymous. One sentence. Submit before you leave.}}

\end{frame}

\begin{frame}{What Were the Key Ideas?}
\note{
% -- NARRATE: What to SAY while showing this slide
``Close notes. 90 seconds. Write 4 most important concepts.''
Walk around the room. Do NOT show next slide yet.

% -- FLEX: [CORE] or [OPTIONAL] + contingency
[CORE] Retrieval practice consolidates learning.
}

\centering
\vspace{1.5cm}
{\Large\bfseries Close your notes.}

\vspace{0.8cm}
{\large Write down the \textbf{4 most important concepts} from today.}

\vspace{0.8cm}
{\normalsize\textcolor{midgray}{90 seconds --- no peeking.}}

\end{frame}

\begin{frame}{Key Takeaways}
\note{
% -- LINK: What prior concept connects to this slide
Students just attempted retrieval. Reveal for self-correction.

% -- NARRATE: What to SAY while showing this slide
Walk through each bullet: ``Serving cost dominates 100x+. Training
optimizes throughput, inference optimizes P99. No universal batching.
KV cache grows with B*S and PagedAttention recovers 60--80\%. Power
of two choices gives exponential improvement. Cold start kills
reactive scaling --- predictive is mandatory.''

% -- FLEX: [CORE] or [OPTIONAL] + contingency
[CORE] Summary after retrieval practice.
IF SHORT: Read bullets quickly.
}

\footnotesize
\begin{itemize}\setlength\itemsep{1pt}
  \item \textbf{Serving cost dominates training cost}: 100$\times$+ for high-volume applications. Every percent of serving efficiency saves significant capital.
  \item \textbf{The fundamental inversion}: Training optimizes throughput (samples/hr). Inference optimizes tail latency (P99 ms). Different metrics, different architectures.
  \item \textbf{No universal batching}: Static for vision, continuous for LLMs (2--4$\times$ throughput), feature-parallel for RecSys. Match technique to workload.
  \item \textbf{KV cache is the bottleneck}: Grows linearly with $B \times S$, often exceeds model weights. PagedAttention recovers 60--80\% of wasted memory.
  \item \textbf{Power of two choices}: Max queue from $O(\log n/\log \log n)$ to $O(\log \log n)$ with one extra probe. Exponential improvement, minimal overhead.
  \item \textbf{Cold start kills reactive scaling}: 5+ min for large models. Predictive scaling and warm pools are mandatory.
\end{itemize}

\end{frame}

\begin{frame}{References}
\note{
% -- NARRATE: What to SAY while showing this slide
``Key papers: vLLM/PagedAttention for KV cache management, Mitzenmacher
for power-of-two-choices, Megatron for tensor parallelism, Leviathan
for speculative decoding.''

% -- FLEX: [OPTIONAL] + contingency
[OPTIONAL] Reference slide for self-study.
}

\footnotesize
\mlsysref{Kwon+23}{Kwon et al. ``PagedAttention for LLM Serving.'' SOSP 2023.}
\mlsysref{Mitzen01}{Mitzenmacher. ``Power of Two Choices in Load Balancing.'' IEEE TPDS, 2001.}
\mlsysref{Shoeybi+19}{Shoeybi et al. ``Megatron-LM: Model Parallelism.'' 2019.}
\mlsysref{Leviathan+23}{Leviathan et al. ``Speculative Decoding.'' ICML 2023.}
\mlsysref{Naumov+19}{Naumov et al. ``DLRM for Recommendation Systems.'' 2019.}
\mlsysref{Yu+22}{Yu et al. ``Orca: Distributed Serving for Generative Models.'' OSDI 2022.}

\end{frame}

\begin{frame}{Next Lecture: Performance Engineering}
\note{
% -- NARRATE: What to SAY while showing this slide
``This chapter built the serving system. Next chapter optimizes every
operation within it: kernel fusion to eliminate memory round-trips,
FlashAttention for O(1) extra memory, and graph compilation to
automate optimizations. Performance engineering turns the serving
system from functional to fast.''

% -- FLEX: [CORE] or [OPTIONAL] + contingency
[CORE] Forward hook plants the question for next lecture.
}

\footnotesize
\begin{columns}[c]
  \begin{column}{0.30\textwidth}
    \centering
    \begin{mlsyscard}{computestroke}
    {\small\bfseries Kernel Fusion}\\[0.05cm]
    {\scriptsize Merge ops, eliminate memory round-trips}
    \end{mlsyscard}
  \end{column}
  \begin{column}{0.30\textwidth}
    \centering
    \begin{mlsyscard}{datastroke}
    {\small\bfseries FlashAttention}\\[0.05cm]
    {\scriptsize Memory-efficient attention: $O(1)$ extra}
    \end{mlsyscard}
  \end{column}
  \begin{column}{0.30\textwidth}
    \centering
    \begin{mlsyscard}{routingstroke}
    {\small\bfseries Compilation}\\[0.05cm]
    {\scriptsize XLA, Triton: graph $\to$ fast kernels}
    \end{mlsyscard}
  \end{column}
\end{columns}

\vspace{0.15cm}
\centering
{\normalsize How do we make every operation in the serving system as fast as possible?}\\[0.05cm]
{\small\textbf{Performance Engineering turns optimization into sustained cost savings.}}

\end{frame}



\appendix

\begin{frame}{Backup: KV Cache Sizing Reference}
\note{
% -- NARRATE: What to SAY while showing this slide
Backup table with KV cache sizes for common models. GQA in Llama-70B
reduces KV heads from 64 to 8, saving 8x cache.

% -- FLEX: [OPTIONAL] Use if students want model-specific numbers.
}

\footnotesize
\renewcommand{\arraystretch}{1.15}
\begin{tabular}{@{}lrrrrr@{}}
  \toprule
  \textbf{Model} & \textbf{L} & \textbf{KV Heads} & \textbf{dim} & \textbf{MB/tok} & \textbf{B=16, S=4K} \\
  \midrule
  Llama-7B   & 32  & 32  & 128 & 0.5  & 32 GB \\
  Llama-70B  & 80  & 8   & 128 & 0.3  & 21 GB \\
  GPT-4 est. & 120 & 128 & 128 & 7.5  & 480 GB \\
  \bottomrule
\end{tabular}

\vspace{0.1cm}
{\footnotesize Grouped Query Attention (GQA) in Llama-70B reduces KV heads from 64 to 8, saving 8$\times$ cache.}

\end{frame}

\begin{frame}{Backup: Continuous vs Static Batching}
\note{
% -- NARRATE: Worked example comparing static and continuous batching.
% -- FLEX: [OPTIONAL] Use for students wanting detailed batching comparison.
}

\footnotesize
\textbf{Example:} 4 requests with generation lengths 10, 50, 20, 100 tokens.

\vspace{0.1cm}
\renewcommand{\arraystretch}{1.15}
\begin{tabular}{@{}lrr@{}}
  \toprule
  & \textbf{Static} & \textbf{Continuous} \\
  \midrule
  Total tokens generated & 400 (4$\times$100 slots) & 180 (actual) \\
  Wasted compute & 55\% & $<$5\% \\
  GPU utilization & 45\% & 95\% \\
  Throughput & 1$\times$ & \textbf{2--4$\times$} \\
  \bottomrule
\end{tabular}

\vspace{0.1cm}
{\footnotesize Static batching pads all requests to the longest generation, wasting 55\% of compute.}

\end{frame}

\begin{frame}{Backup: Speculative Decoding Math}
\note{
% -- NARRATE: Acceptance probability and expected speedup for speculative decoding.
% -- FLEX: [OPTIONAL] Use when students ask about the speculative decoding acceptance rate.
}

\footnotesize
\textbf{Speculative Decoding Speedup Formula:}

\vspace{0.1cm}
Draft model generates $K$ candidate tokens. Target model verifies all $K$ in one forward pass.

\vspace{0.1cm}
{\small
$$\text{Speedup} = \frac{K + 1}{1 + (1 - p_{\text{accept}}) \cdot K \cdot c}$$
}

\vspace{0.05cm}
\renewcommand{\arraystretch}{1.1}
\begin{tabular}{@{}lrrr@{}}
  \toprule
  \textbf{Acceptance Rate} & \textbf{K=4} & \textbf{K=8} & \textbf{Effective Speedup} \\
  \midrule
  90\% (easy text) & 3.6 tokens/step & 7.2 tokens/step & 2.5--3$\times$ \\
  70\% (general)    & 2.8 tokens/step & 5.6 tokens/step & 1.8--2.2$\times$ \\
  50\% (creative)   & 2.0 tokens/step & 4.0 tokens/step & 1.3--1.5$\times$ \\
  \bottomrule
\end{tabular}

\vspace{0.1cm}
{\footnotesize\textbf{Key insight:} Speedup depends on draft--target agreement. Easy tasks (code completion) see higher acceptance than creative tasks (story writing).}

\end{frame}


\end{document}
